\documentclass[11pt]{article}
\usepackage[T1]{fontenc}
\usepackage{lmodern}
\usepackage[margin=1in]{geometry}
\usepackage{amsmath,amssymb,amsthm}
\usepackage{microtype}
\usepackage[colorlinks=true,linkcolor=blue,citecolor=blue,urlcolor=blue]{hyperref}
\hypersetup{pdftitle={Quadratic barrier parameters for hyperbolic balls},pdfauthor={Kaiwen Jin}}
\newtheorem{theorem}{Theorem}[section]
\newtheorem{lemma}[theorem]{Lemma}
\newtheorem{proposition}[theorem]{Proposition}
\newcommand{\HH}{\mathbb H}
\newcommand{\R}{\mathbb R}
\newcommand{\SC}{\nabla^3}
\title{Quadratic barrier parameters for hyperbolic balls}
\author{Kaiwen Jin\\\small\href{mailto:yui.ui.siki@gmail.com}{yui.ui.siki@gmail.com}}
\date{October 14, 2026}

\begin{document}
\maketitle
\begin{abstract}
For geodesic balls of radius \(R\) in the hyperbolic plane of curvature
\(-1\), we prove that the least parameter of a self-concordant barrier is
of order \(R^2\). Self-concordance is required for arbitrary triples of
tangent vectors. The lower bound holds for every barrier, without a
symmetry assumption: at a maximum on a geodesic circle, curvature forces
a fixed positive radial Hessian. A scalar convexity argument turns this
local constraint into a quadratic lower bound. The known quadratic
construction is therefore optimal in order, and an \(O(\sqrt R)\)
parameter is impossible for this definition.
\end{abstract}

\section{The parameter of a hyperbolic ball}

Self-concordant barriers extend interior-point methods to geodesically
convex optimization on manifolds. Their barrier parameter enters the
iteration bounds, so its dependence on the size of the domain is a basic
question. For hyperbolic balls, constructions with a quadratic dependence
on the radius are known \cite{Hirai,HNW}. Criscitiello, Nieuwboer and Walter
proved a lower bound of order \(\sqrt R\), and asked whether it is sharp
\cite[Corollary~1.4 and the following paragraph]{CNW}. Their lower bound
also applies to the weaker notion of self-concordance along geodesics.
Here we determine the order for the full mixed-derivative definition.

Let \(\HH^2\) have curvature \(-1\), fix \(o\in\HH^2\), and write
\[
 D_R=\{x\in\HH^2:d(o,x)<R\}.
\]
All derivatives use the Levi-Civita connection. For \(F\in C^3(D_R)\),
put \(H=\nabla^2F\) and
\[
 T(u,v,w)=(\nabla_uH)(v,w).
\]
For \(\theta\geq0\), a \(\theta\)-barrier means that \(H\) is positive
definite, the epigraph of \(F\) is closed in \(\HH^2\times\R\), and, at
every \(x\in D_R\) and for all \(u,v,w\in T_x\HH^2\),
\begin{align}
 |T(u,v,w)|&\leq2\sqrt{H(u,u)H(v,v)H(w,w)},\label{eq:SC}\\
 |dF(v)|^2&\leq\theta H(v,v).\label{eq:gradient}
\end{align}
This is the full definition in \cite[Definition~1.1 and Definition~4.1]{HNW}
and \cite[Section~2.3]{CNW}. Define
\[
 \theta_*(R)=\inf\{\theta\geq0:
       \text{there is a \(\theta\)-barrier on \(D_R\)}\}.
\]
The infimum allows a different function for each \(R\).

\begin{theorem}\label{thm:main}
For every \(R>1\),
\[
 \frac{(R-1)^2}{256}\leq\theta_*(R)\leq1+2R^2.
\]
In particular, for every \(R\geq2\),
\[
 \frac{R^2}{1024}\leq\theta_*(R)\leq\frac94R^2.
\]
Thus \(\theta_*(R)\) has order \(R^2\) as \(R\to\infty\), with absolute
comparison constants, and is not \(O(\sqrt R)\).
\end{theorem}

The upper bound is a specialization of the construction in
\cite[Theorem~1.6]{HNW}. The new lower bound uses the asymmetry of the
third covariant derivative. At a circle maximum, the tangential Hessian
is controlled by the outward derivative; curvature then forces a
positive radial second derivative. Taking maxima over circles preserves
enough scalar convexity to compare this growth with
\eqref{eq:gradient}. No averaging or rotational symmetry of \(F\) is
used. The argument does not give the same lower bound when
\eqref{eq:SC} is assumed only for \(u=v=w\); the corresponding weaker
question remains separate.

\section{The lower bound}

For \(r=d(o,x)>0\), use geodesic polar coordinates. The metric and
orthonormal radial and tangential vectors are
\[
 g=dr^2+\sinh^2r\,d\phi^2,\qquad
 e_r=\partial_r,\qquad e_t=\frac1{\sinh r}\partial_\phi.
\]
In this frame, \(\nabla_{e_t}e_t=-\coth(r)e_r\). We use the curvature
convention
\[
 {\cal R}(U,V)W=\nabla_U\nabla_VW-\nabla_V\nabla_UW-\nabla_{[U,V]}W.
\]
Commuting the two covariant derivatives of the covector \(dF\), and using
torsion-freeness, gives
\begin{equation}
 T(u,v,w)-T(v,u,w)=-dF({\cal R}(u,v)w).\label{eq:commutator}
\end{equation}
The minus sign is the curvature action on a covector. Since the curvature
is \(-1\), \({\cal R}(e_r,e_t)e_t=-e_r\).

\begin{lemma}\label{lem:circle}
Suppose \(H>0\) and \eqref{eq:SC} holds. At every maximum of \(F\) on a
circle of radius \(r\in(0,R)\),
\[
 F_{rr}\geq\frac{\tanh^2r}{16}.
\]
\end{lemma}

\begin{proof}
At such a point let \(p=F_r\), \(a=H(e_r,e_r)\), and \(b=H(e_t,e_t)\).
The second derivative in the angular variable is nonpositive. Therefore
\begin{equation}
 0<b=\frac{F_{\phi\phi}}{\sinh^2r}+p\coth r\leq p\coth r,
 \label{eq:tangential}
\end{equation}
which proves \(p>0\). By \eqref{eq:commutator} and the two mixed instances
of \eqref{eq:SC},
\[
 p=|T(e_r,e_t,e_t)-T(e_t,e_r,e_t)|
   \leq4\sqrt a\,b
   \leq4\sqrt a\,p\coth r.
\]
Dividing by \(p>0\) proves the assertion, since \(a=F_{rr}\).
\end{proof}

\begin{proposition}\label{prop:lower}
Every \(\theta\)-barrier on \(D_R\), with \(R>1\), satisfies
\(\theta\geq(R-1)^2/256\).
\end{proposition}

\begin{proof}
Let
\[
 m(r)=\max_{\phi}F(r,\phi),\qquad 0\leq r<R.
\]
Each maximum exists because the circle is compact. On every compact
radial interval inside the ball, the derivatives of the radial
restrictions are bounded uniformly in the direction. The mean value
theorem gives the same Lipschitz bound for their maximum.

Extend \(m\) to \((-R,R)\) as the maximum of
\(F(\exp_o(r\xi))\) over unit vectors \(\xi\in T_o\HH^2\). Each restriction
is convex, so the maximum is convex. It is even, since replacing \(r\)
by \(-r\) replaces \(\xi\) by \(-\xi\). It follows that \(m\) is
nondecreasing for \(r\geq0\).

We claim that
\begin{equation}
 h(r)=m(r)-\frac{r^2}{128}
       \quad\hbox{is convex on }[1,R).\label{eq:strong}
\end{equation}
Indeed, if chord convexity fails on some \([s,t]\subset[1,R)\), let
\(L\) be the affine function interpolating \(h(s)\) and \(h(t)\).
The function
\[
 F(r,\phi)-\frac{r^2}{128}-L(r)
\]
has a positive maximum on the compact annulus \([s,t]\times S^1\).
Its maximum on either endpoint circle is zero, so a maximizing point
\((r_0,\phi_0)\) has \(s<r_0<t\). Its angle also maximizes
\(F(r_0,\cdot)\). The radial second derivative test gives
\(F_{rr}(r_0,\phi_0)\leq1/64\), whereas Lemma~\ref{lem:circle} gives
\[
 F_{rr}(r_0,\phi_0)\geq\frac{\tanh^2r_0}{16}>
 \frac1{64}.
\]
Here \(r_0\geq1\) and \(\tanh1>1/2\). This proves \eqref{eq:strong},
including when the maximizing angle changes or \(m\) is not
differentiable.

Finite locally Lipschitz convex functions have finite right derivatives,
which are nondecreasing. Applying this to \(h\), and using
\(m'_+(1)\geq0\), yields
\begin{equation}
 m'_+(r)\geq\frac{r-1}{64},\qquad 1<r<R.\label{eq:lower-slope}
\end{equation}

For the opposing estimate, first note that \(\theta>0\): otherwise
\eqref{eq:gradient} makes \(dF\) vanish everywhere, contradicting \(H>0\).
For any unit \(\xi\in T_o\HH^2\), put
\(f(t)=F(\exp_o(t\xi))\), \(-R<t<R\). Along this geodesic,
\[
 f''(t)>0,\qquad f'(t)^2\leq\theta f''(t).
\]
If \(f'(s)>0\), then \(f'(t)>0\) for \(s<t<R\), and
\[
 \left(\frac1{f'}\right)'=-\frac{f''}{(f')^2}\leq-\frac1\theta.
\]
Comparison on \([s,t]\), with \(t<R\), gives
\[
 0<\frac1{f'(t)}
   \leq\frac1{f'(s)}-\frac{t-s}{\theta}.
\]
Letting \(t\uparrow R\) proves \(f'(s)\leq\theta/(R-s)\). The same
inequality holds if \(f'(s)\leq0\). Thus, for every direction and
\(0\leq s<t<R\),
\[
 f(t)-f(s)\leq\theta\log\frac{R-s}{R-t}.
\]
This follows, for example, by differentiating
\(f(t)+\theta\log(R-t)\), whose derivative is nonpositive. Taking
maxima over all directions gives the identical increment bound for
\(m\), and therefore
\begin{equation}
 m'_+(s)\leq\frac{\theta}{R-s},\qquad 0\leq s<R.\label{eq:upper-slope}
\end{equation}
No interchange between differentiation and a maximum is needed.

Combining \eqref{eq:lower-slope} and \eqref{eq:upper-slope} at
\(s=(R+1)/2\) gives
\[
 \theta\geq\frac{(s-1)(R-s)}{64}=\frac{(R-1)^2}{256}.
\]
\end{proof}

\section{The upper bound}

We include the specialization of the known construction to check every
clause of the barrier definition.

\begin{proposition}\label{prop:upper}
For every \(R>0\), the function
\[
 F_R(x)=-\log\bigl(R^2-d(o,x)^2\bigr)+d(o,x)^2
\]
is a \((1+2R^2)\)-barrier on \(D_R\).
\end{proposition}

\begin{proof}
Put \(q(x)=d(o,x)^2\), \(Q=\nabla^2q\).
In the source convention, \(\alpha\)-self-concordance replaces the
coefficient \(2\) in \eqref{eq:SC} by \(2/\sqrt\alpha\).
By \cite[Theorem~5.16(i)]{HNW}, \(q\) is \(8\)-self-concordant on
curvature \(-1\) hyperbolic space, hence also \(1\)-self-concordant.
This assertion uses the full tensor definition. To make the repeated-slot
form of the source estimate explicit, fix the first vector and regard the
last two slots as a symmetric bilinear form \(A\). If
\(|A(z,z)|\leq KQ(z,z)\) for all \(z\), polarization gives, for \(a>0\),
\[
 |A(v,w)|\leq\frac K2\bigl(a^2Q(v,v)+a^{-2}Q(w,w)\bigr).
\]
For nonzero \(v,w\), choosing
\(a^2=\sqrt{Q(w,w)/Q(v,v)}\) gives
\(|A(v,w)|\leq K\sqrt{Q(v,v)Q(w,w)}\); zero vectors are immediate.
Thus no constant is lost in obtaining arbitrary triples.

The epigraph construction \cite[Theorem~1.6]{HNW}, applied with
\(\alpha=1\), states that
\[
 B(x,S)=-\log(S-q(x))+q(x),\qquad S>q(x),
\]
on \(\HH^2\times\R\) is strongly \(1\)-self-concordant and satisfies
\[
 |dB(V)|^2\leq(1+2q(x))(\nabla^2B)(V,V)
\]
for every product tangent vector \(V\). Restrict to \(S=R^2\).
The product connection obeys
\[
 \nabla^{\HH^2\times\R}_{(u,0)}(v,0)=(\nabla^{\HH^2}_uv,0),
\]
so both covariant derivatives in the definition pull back exactly.
The restriction is \(F_R\), and \(q<R^2\) gives
\eqref{eq:gradient} with \(\theta=1+2R^2\).

Squared distance is smooth also at \(o\). With \(\delta=R^2-q>0\),
\[
 \nabla^2F_R=(1+\delta^{-1})Q+\delta^{-2}dq\otimes dq.
\]
The eigenvalues of \(Q\) are \(2\) and \(2r\coth r\) for \(r>0\),
and \(Q=2g\) at \(o\); hence this Hessian is positive definite.
Finally \(F_R\to+\infty\) as \(r\uparrow R\).
If an epigraph sequence converges to a finite point of
\(\HH^2\times\R\), this divergence excludes a limit on \(r=R\);
an interior limit lies in the epigraph by continuity. Thus the epigraph
is closed in the ambient space.
\end{proof}

\begin{proof}[Proof of Theorem~\ref{thm:main}]
Proposition~\ref{prop:lower} applies to every admissible barrier, so its
lower bound survives the infimum defining \(\theta_*(R)\).
Proposition~\ref{prop:upper} supplies the upper bound.
For \(R\geq2\), use \(R-1\geq R/2\) and \(1\leq R^2/4\).
The resulting positive quadratic lower bound also excludes
\(\theta_*(R)=O(\sqrt R)\).
\end{proof}

\paragraph{AI use.}
OpenAI Codex assisted in developing and checking the argument and in
preparing this manuscript.

\begin{thebibliography}{9}
\bibitem{CNW}
C. Criscitiello, H. Nieuwboer and M. Walter,
\emph{Negative curvature obstructs the existence of good barriers for
interior-point methods},
arXiv:2604.07319v1 (2026).
\href{https://arxiv.org/abs/2604.07319v1}{arXiv:2604.07319v1}.

\bibitem{HNW}
H. Hirai, H. Nieuwboer and M. Walter,
\emph{Interior-Point Methods on Manifolds: Theory and Applications},
Foundations of Computational Mathematics (2026).
\href{https://doi.org/10.1007/s10208-026-09756-8}{doi:10.1007/s10208-026-09756-8}.
Also \href{https://arxiv.org/abs/2303.04771v2}{arXiv:2303.04771v2} (2023).

\bibitem{Hirai}
H. Hirai,
\emph{On a manifold formulation of self-concordant functions},
arXiv:2212.10981v2 (2023).
\href{https://arxiv.org/abs/2212.10981v2}{arXiv:2212.10981v2}.
\end{thebibliography}
\end{document}
